% Nota de figura com argumento (não existe no projeto base)
\newcommand{\notafigv}[1]{\par\noindent{\footnotesize\textit{Nota:} #1\newline Fonte: elaboração do autor.}}

\chapter{Resultados: Ideologia e Formação Econômica como Preditores Diferenciais}
\label{capitulo_4}

Os resultados revelam dois padrões associativos sobrepostos. No primeiro, o espectro político autodeclarado é o preditor mais consistente, mais robusto e mais ubíquo das 53 variáveis dependentes analisadas. No segundo, a formação econômica avançada está associada a respostas sistematicamente distintas em pontos específicos — precisamente aqueles onde o item tem componente factual mais dominante que identitária. Dois padrões, portanto, mas com escalas e abrangências distintas.

Esse resultado é compatível com o que \citeonline{Sapienza2013} documentaram comparando economistas especialistas e americanos médios: tanto a orientação política quanto a formação técnica estão associadas a divergências de crenças, com padrões diferenciados conforme a natureza do item. O que se acrescenta aqui é a dimensão quantitativa desse fenômeno no contexto brasileiro e a constatação de que, quando os dois preditores coexistem, eles não se cancelam — a posição ideológica prediz a maior parte das respostas e a formação econômica prediz um subconjunto específico do resíduo não explicado pela ideologia.

Este capítulo está organizado em quatro movimentos. Primeiro (\autoref{sec:perfil-amostral} e \autoref{sec:visao-panoramica}), descreve-se o padrão global: há estrutura, não ruído. Segundo (\autoref{sec:dominancia-espectro}), demonstra-se que o padrão é majoritariamente ideológico. Terceiro (\autoref{sec:formacao-economica} e \autoref{sec:itens-chave}), examina-se onde e como a formação econômica atua. Quarto (\autoref{sec:bh-virtude} e \autoref{sec:sintese-hipoteses}), extraem-se implicações sem ultrapassar o que os dados permitem.

O questionário ficou aberto de 1.º de agosto de 2025 a 29 de setembro de 2025, obtendo $n = 184$ respostas válidas (91\% público geral, $n \approx 167$; 9\% economistas com pós-graduação, $n = 17$). Dos 184 respondentes, 38 não declararam posição ideológica (``Independente'' ou ``Sem opinião'') e recebem \texttt{espectro = NaN}, sendo excluídos dos modelos que incluem o controle de espectro político — o que corresponde a todos os 53 modelos. O N efetivo por modelo situa-se entre 144 e 146. Foram estimados 53 modelos de regressão ordinal — 48 pelo método de chances proporcionais (PO) e 5 pelo método Firth-Ridge para DVs com Hessiana singular —, com controle por espectro político, escolaridade, faixa etária, gênero e exposição profissional à economia. O perfil descritivo da amostra consta de \autoref{apendice:tab_descritiva_v2}.

\section{Perfil Amostral}
\label{sec:perfil-amostral}

As estatísticas descritivas completas por DV e por grupo (N total, médias e desvios-padrão para \texttt{econ}=0 e \texttt{econ}=1) constam de \autoref{apendice:tab_descritiva_v2}. Em resumo: o grupo de controle inclui $n \approx 167$ respondentes e o grupo de economistas $n = 17$, com N efetivo por modelo entre 144 e 146 após exclusão de respondentes que não informaram posição ideológica (\texttt{espectro = NaN}; $n_{\text{excl}} = 38$). A distribuição por espectro político nos dois grupos é apresentada na Tabela~\ref{tab:crosstab_econ_esp} (Seção~\ref{sec:tamanho_amostral}).

\section{Visão Panorâmica: O Mapa do Argumento}
\label{sec:visao-panoramica}

\begin{figure}[htbp]
\centering
\caption{Mapa de calor dos efeitos por bloco temático: espectro político (BH-corrigido) e formação econômica (nominal).}
\includegraphics[width=0.95\textwidth]{Textuais/analise_v2/outputs/figures/heatmaps/bloc_heatmap.png}
\notafigv{Cada célula indica a significância estatística do coeficiente no modelo PO ou Firth-Ridge. Azul = espectro político BH-significativo; vermelho = \texttt{econ} nominalmente significativo ($p < 0{,}05$, sem correção BH). Os 53 itens estão organizados pelos seis blocos temáticos pré-especificados.}
\label{fig:bloc_heatmap}
\end{figure}

Antes de qualquer tabela de coeficientes, a Figura~\ref{fig:bloc_heatmap} oferece o mapa geral do argumento. Ela mostra dois padrões sobrepostos que sintetizam todo o capítulo. O espectro político (azul) aparece de modo amplo e sistemático: coloração densa em praticamente todos os blocos, com intensidade máxima em H3 (Antimercado) e H7 (Ideologia). A formação econômica (vermelho) aparece em manchas pontuais: quando o tema é politicamente carregado, o espectro domina; quando o tema tem componente factual mais dominante, a formação econômica exibe diferenças condicionais pontuais.

Esse panorama antecipa o resultado central: o espectro político é o preditor dominante e persistente das percepções econômicas nesta amostra; a formação econômica prediz diferenças em pontos específicos, sem efeito discernível nos demais. A Tabela~\ref{tab:sintese_resumo} apresenta uma síntese compacta dos resultados por bloco, com os 20 itens mais informativos do conjunto.


% Tabela gerada automaticamente por latex_export.py
\begingroup
\footnotesize
\SingleSpacing
\begin{longtable}{%
  >{\RaggedRight\arraybackslash}p{4.5cm}%
  >{\centering\arraybackslash}p{1.4cm}%
  >{\centering\arraybackslash}p{1.4cm}%
  >{\centering\arraybackslash}p{1.2cm}%
  >{\centering\arraybackslash}p{1.4cm}%
  >{\centering\arraybackslash}p{1.4cm}%
  >{\centering\arraybackslash}p{1.2cm}}
\caption{Resultados selecionados --- 20 DVs com menor $p$ nominal ($\alpha=0,05$; 53 testes)}\label{tab:sintese_resumo}\\
\toprule
\textbf{DV} & $\hat{\beta}_\text{econ}$ & $p$ (nom.) & $p$ (BH) & $\hat{\beta}_\text{esp}$ & $p_\text{esp}$ (nom.) & $R^2_\text{McF.}$ \\
\midrule
\endfirsthead
\multicolumn{7}{r}{\footnotesize \textit{(continuação)}}\\\toprule
\textbf{DV} & $\hat{\beta}_\text{econ}$ & $p$ (nom.) & $p$ (BH) & $\hat{\beta}_\text{esp}$ & $p_\text{esp}$ (nom.) & $R^2_\text{McF.}$ \\
\midrule
\endhead
\midrule
\multicolumn{7}{r}{\footnotesize \textit{(continua)}}\\
\endfoot
\bottomrule
\endlastfoot
  déficit\_\allowbreak{}federal\_\allowbreak{}grande & -1,599 & 0,014* & 0,364 & 1,068 & $<$ 0,001*** & 0,182 \\
  produtividade\_\allowbreak{}aumentando\_\allowbreak{}devagar & 1,980 & 0,017* & 0,364 & 0,087 & 0,581 & 0,114 \\
  reforma\_\allowbreak{}previdência\_\allowbreak{}necessária & -1,395 & 0,021* & 0,364 & 0,652 & $<$ 0,001*** & 0,084 \\
  empresas\_\allowbreak{}estão\_\allowbreak{}enviando & 1,390 & 0,033* & 0,403 & 0,094 & 0,637 & 0,048 \\
  produtos\_\allowbreak{}importados\_\allowbreak{}benéficos & 1,449 & 0,039* & 0,403 & 0,817 & $<$ 0,001*** & 0,123 \\
  reforma\_\allowbreak{}trabalhista\_\allowbreak{}necessária & -1,179 & 0,052 & 0,403 & 0,626 & $<$ 0,001*** & 0,074 \\
  taxa\_\allowbreak{}juros\_\allowbreak{}selic & 1,616 & 0,053 & 0,403 & -0,402 & 0,020* & 0,099 \\
  empresas\_\allowbreak{}estão\_\allowbreak{}reduzindo & 1,086 & 0,074 & 0,480 & -0,563 & $<$ 0,001*** & 0,078 \\
  governo\_\allowbreak{}atual\_\allowbreak{}sabe & 1,149 & 0,082 & 0,480 & -1,978 & $<$ 0,001*** & 0,369 \\
  redução\_\allowbreak{}recente\_\allowbreak{}postos & -1,390 & 0,092 & 0,488 & 0,503 & 0,006** & 0,069 \\
  últimos\_\allowbreak{}20\_\allowbreak{}anos & 1,018 & 0,106 & 0,511 & 0,349 & 0,054 & 0,040 \\
  educação\_\allowbreak{}qualificação\_\allowbreak{}profissional & 0,999 & 0,132 & 0,585 & 0,233 & 0,146 & 0,043 \\
  brasil\_\allowbreak{}chance\_\allowbreak{}virar & -0,845 & 0,147 & 0,600 & -0,569 & 0,001** & 0,055 \\
  entrada\_\allowbreak{}estrangeiros\_\allowbreak{}mercado & 0,825 & 0,181 & 0,687 & -0,022 & 0,893 & 0,023 \\
  tecnologia\_\allowbreak{}causa\_\allowbreak{}demissões & 0,733 & 0,216 & 0,716 & -0,193 & 0,301 & 0,083 \\
  deduções\_\allowbreak{}demais\_\allowbreak{}empresas & 0,683 & 0,216 & 0,716 & 0,391 & 0,014* & 0,028 \\
  competição\_\allowbreak{}entre\_\allowbreak{}empresas & 1,182 & 0,294 & 0,788 & 0,867 & $<$ 0,001*** & 0,127 \\
  aumento\_\allowbreak{}uso\_\allowbreak{}tecnologia$\dag$ & 1,991 & 0,303 & 0,788 & 0,955 & 0,038* & 0,219 \\
  acordos\_\allowbreak{}comerciais\_\allowbreak{}outros$\dag$ & -1,276 & 0,316 & 0,788 & -0,584 & 0,207 & 0,154 \\
  corrupção\_\allowbreak{}principal\_\allowbreak{}causa & -0,576 & 0,330 & 0,788 & 1,002 & $<$ 0,001*** & 0,150 \\
\bottomrule
\end{longtable}
\par\smallskip\noindent\begin{minipage}{\linewidth}\footnotesize\sloppy\centering
Fonte: elaboração do autor.\\
\textit{\textbf{Notas:}}\ $\dag$ = Hessiana singular (Firth-Ridge); *\,$p<0,05$, **\,$p<0,01$, ***\,$p<0,001$. BH = Benjamini-Hochberg FDR. Ordenado por $p$ nominal crescente; 53 DVs no total.
\end{minipage}
\endgroup


\section{A Dominância do Espectro Político}
\label{sec:dominancia-espectro}

\begin{figure}[htbp]
\centering
\caption{\textit{Forest plot} dos coeficientes de espectro político ($\hat{\beta}_{\text{esp}}$) por DV.}
\includegraphics[width=0.95\textwidth]{Textuais/analise_v2/outputs/figures/forestplots/espectro_forest.png}
\notafigv{Cada ponto é o coeficiente estimado de \texttt{espectro} no modelo PO (ou Firth-Ridge); barras de erro = IC95\% assintótico. Pontos coloridos = BH-significativos; pontos cinza = não significativos após BH. DVs organizadas por bloco temático. Coeficiente positivo indica que respondentes mais à direita concordam mais com a afirmação.}
\label{fig:espectro_forest}
\end{figure}

O espectro político apresenta efeito estatisticamente significativo após correção de Benjamini-Hochberg em \textbf{34 das 53 DVs} (64\%), distribuídos de modo heterogêneo mas sistematicamente denso pelos blocos temáticos: H3-Antimercado alcança 9/9 (100\%); H7-Ideologia atinge 11/12 (92\%); H4-Antiestrangeiro, H5-Antitrabalho e H6-Pessimista apresentam taxas entre 40\% e 67\%. Apenas o bloco de Percepção Específica exibe menor densidade (2/3). A Figura~\ref{fig:espectro_forest} torna visível a estrutura: coeficientes positivos onde a direita concorda mais (déficit, reformas, produtividade), negativos onde a esquerda concorda mais (lucros excessivos, regulação, desigualdade). A magnitude típica é $|\hat{\beta}_{\text{esp}}| \in [0{,}4,\; 0{,}9]$ em log-odds.

\paragraph{O filtro ideológico como mecanismo, não apenas como preditor.} O resultado não deve ser lido como simples correlação entre posição e crença. \citeonline{kahan2012polarization} chama esse fenômeno de \textit{cognição cultural}: a predisposição ideológica não apenas colore a interpretação de evidências ambíguas — ela determina quais afirmações sequer chegam a ser ponderadas de forma imparcial. Em termos de \citeonline{taber2006motivated}, o cidadão avalia argumentos congruentes com suas crenças prévias com menor ceticismo e argumentos incongruentes com ceticismo seletivo. Como em Cervantes, o problema não está apenas no objeto visto, mas na narrativa que o antecede: o moinho pode estar diante dos olhos, mas a imaginação já o converteu em gigante. Isso explica por que o efeito aparece tanto em questões factuais (déficit, produtividade) quanto em questões de valor (reformas, desigualdade): o filtro ideológico opera antes da separação entre fato e julgamento.

\citeonline{zaller1992nature} oferece a moldura sistematizada: a opinião pública é recepção e aceitação seletiva de mensagens. Cidadãos aceitam mais facilmente mensagens que reforçam suas predisposições e resistem a mensagens que as contradizem, mesmo quando a informação é idêntica. A consequência direta para este estudo é que o preditor dominante observado ($\hat{\beta}_{\text{esp}}$) captura a predisposição ideológica — estruturalmente anterior à formação econômica no processo de formação de crenças — e explica a maior parte da variância nas 53 DVs.

Uma advertência metodológica é necessária: causalidade reversa é plausível em parte dos itens. Em H7 especialmente (reformas, avaliação do governo, taxa Selic), economistas com formação avançada podem ter desenvolvido posições ideológicas influenciadas pela própria formação — o que não invalida o resultado, mas convida à prudência interpretativa. O preditor dominante observado é correlacional, não estrutural.

\section{Formação Econômica: Sinal Substantivo sob Baixa Potência}
\label{sec:formacao-economica}

\begin{figure}[htbp]
\centering
\caption{\textit{Forest plot} dos coeficientes de formação econômica ($\hat{\beta}_{\texttt{econ}}$) com intervalos de confiança bootstrap BC$_0$.}
\includegraphics[width=0.95\textwidth]{Textuais/analise_v2/outputs/figures/forestplots/econ_forest.png}
\notafigv{Pontos coloridos = nominalmente significativos ($p < 0{,}05$), sem correção BH; pontos cinza = não significativos. IC95\% bootstrap BC$_0$ ($B=999$) para os quatro itens-chave; ICs assintóticos para os demais. Coeficiente negativo indica que economistas discordam mais da afirmação.}
\label{fig:econ_forest}
\end{figure}

A formação econômica avançada apresenta \textbf{5 efeitos nominalmente significativos} ($p < 0{,}05$) sobre as 53 DVs. Após correção de Benjamini-Hochberg, \textbf{nenhum desses efeitos sobrevive}: o menor $p_{\text{BH}}$ é $0{,}364$. A leitura correta desse resultado não é "formação econômica não importa" — é que a amostra não tem potência para confirmar o que os sinais sugerem.

Com $n_{\texttt{econ}} = 17$, o poder estatístico para detectar OR $= 2{,}0$ é apenas $28{,}2\%$ (\autoref{apendice:tab_power_v2}). Isso significa que, mesmo na presença de efeitos reais de magnitude moderada, este estudo teria probabilidade de $71{,}8\%$ de não rejeitá-los. Há sinais substantivos coerentes com a teoria; a amostra simplesmente não autoriza conclusão confirmatória ampla. Declarar isso explicitamente é virtude metodológica, não desculpa.

O que os 5 efeitos nominais comunicam é a \textit{direção} do padrão associativo. Todos apontam na direção esperada segundo a literatura de consenso econômico: respondentes com pós-graduação em Economia exibem menor concordância com afirmações populistas sobre déficit e importações; exibem maior concordância com a preocupação sobre produtividade — em linha com \citeonline{Fuller2014} para o consenso econômico americano e com os vieses catalogados por \citeonline{The_Myth_of_the_Rational_Voter}. O retrodesign de \citeonline{Gelman2014} acrescenta que as magnitudes nominais estão provavelmente infladas por fator de $1{,}2$--$1{,}4\times$ (Erro Tipo M): os efeitos reais são menores, não maiores, do que os observados. Em síntese: o estudo tem poder exploratório, não confirmatório, e os 5 sinais nominais são evidência compatível com associações de magnitude moderada — indício que justifica estudo confirmatório com $n_{\texttt{econ}} \geq 120$.

\section{A Correção por Multiplicidade como Virtude Epistemológica}
\label{sec:bh-virtude}

A correção de Benjamini-Hochberg \cite{benjamini1995fdr} não é mero procedimento estatístico — é o equivalente empírico de uma norma de autocrítica. Com 53 hipóteses testadas simultaneamente, a probabilidade de encontrar pelo menos um falso positivo ao nível $\alpha = 0{,}05$ sem correção seria de $1 - (1-0{,}05)^{53} \approx 93\%$. A correção BH controla a taxa de falsas descobertas: ela funciona como freio contra a tentação de transformar ruído em descoberta.

A ironia metodológica é precisa: este trabalho diagnostica crenças econômicas confortáveis do eleitorado. Sem a correção por multiplicidade, o pesquisador correria o risco de confirmar sua própria hipótese preferida pela mesma lógica que critica — escolhendo, entre muitos testes, os que sustentam a narrativa desejada. \citeonline{benjamini1995fdr} fornece o mecanismo; \citeonline{popperlogic} fornece a motivação filosófica: a ciência avança pela disposição de refutar, não de confirmar. Aplicar BH é, nesse sentido, tratar a si mesmo com a mesma severidade que se aplica ao objeto de estudo.

\section{Itens-Chave: Magnitude dos Efeitos e Público Contrafactual}
\label{sec:itens-chave}

Quatro itens foram selecionados a priori como ilustrativos da diversidade de mecanismos: dois com coeficiente negativo de \texttt{econ} (economistas mais céticos com afirmações populistas), um com coeficiente positivo (economistas mais preocupados com problema estrutural pouco saliente) e um que combina abertura comercial com viés antiestrangeiro.


% Tabela gerada automaticamente por latex_export.py
\begingroup
\footnotesize
\SingleSpacing
\begin{longtable}{%
  >{\RaggedRight\arraybackslash}p{4.5cm}%
  >{\centering\arraybackslash}p{1.4cm}%
  >{\centering\arraybackslash}p{1.4cm}%
  >{\centering\arraybackslash}p{1.4cm}%
  >{\centering\arraybackslash}p{1.4cm}%
  >{\centering\arraybackslash}p{1.4cm}%
  >{\centering\arraybackslash}p{1.4cm}%
}
\caption{Efeitos Marginais Médios (AME) de \texttt{econ} para os itens-chave}\label{tab:ame_chave}\\
\toprule
\textbf{DV} & $\bar{Y}_{e=0}$ & $\bar{Y}_{e=1}$ & $\Delta \bar{Y}$ & AME cat0 & AME cat1 & AME cat2 \\
\midrule
\endfirsthead
\multicolumn{7}{r}{\footnotesize \textit{(continuação)}}\\\toprule
\textbf{DV} & $\bar{Y}_{e=0}$ & $\bar{Y}_{e=1}$ & $\Delta \bar{Y}$ & AME cat0 & AME cat1 & AME cat2 \\
\midrule
\endhead
\midrule
\multicolumn{7}{r}{\footnotesize \textit{(continua)}}\\
\endfoot
\bottomrule
\endlastfoot
  déficit\_\allowbreak{}federal\_\allowbreak{}grande & 1,476 & 1,016 & -0,460 & 0,169 & 0,123 & -0,291 \\
  gasto\_\allowbreak{}ajuda\_\allowbreak{}externa & 0,775 & 0,639 & -0,136 & 0,085 & -0,035 & -0,050 \\
  produtividade\_\allowbreak{}aumentando\_\allowbreak{}devagar & 1,120 & 1,722 & 0,601 & -0,201 & -0,200 & 0,401 \\
  produtos\_\allowbreak{}importados\_\allowbreak{}benéficos & 1,263 & 1,628 & 0,366 & -0,069 & -0,229 & 0,297 \\
\bottomrule
\end{longtable}
\par\smallskip\noindent\begin{minipage}{\linewidth}\footnotesize\sloppy\centering
Fonte: elaboração do autor.\\
\textit{\textbf{Notas:}}\ AME via diferenças finitas. $\Delta\bar{Y} = \sum_j j \cdot \widehat{\text{AME}}_j$. Categorias em escala 0--2. 4 DVs reportados.
\end{minipage}
\endgroup


A Tabela~\ref{tab:ame_chave} quantifica a magnitude das diferenças condicionais. O $\Delta\bar{Y}$ indica a diferença esperada na escala 0--2 entre economistas e o grupo de controle, após condicionar nas demais covariáveis. Os valores nominais são: $-0{,}46$ para déficit, $-0{,}14$ para ajuda externa, $+0{,}60$ para produtividade, $+0{,}37$ para importados. Esses valores devem ser interpretados com duas ressalvas. Primeira: expressam associações condicionais, não efeitos causais da formação econômica. Segunda: dado o poder estatístico de $\approx 15$--$28\%$ para $n_{\texttt{econ}} = 17$ (dependendo do item), o retrodesign de \citeonline{Gelman2014} indica razão de exagero esperada entre $1{,}5\times$ e $2{,}5\times$ para as estimativas que atingiram significância nominal; os valores pontuais reportados são, portanto, limites superiores plausíveis da magnitude real. A direção dos efeitos (sinal de $\hat{\beta}_{\texttt{econ}}$) é mais informativa do que as magnitudes absolutas.

As figuras a seguir traduzem essas diferenças no trio público--predição-contrafactual--economistas. O "contrafactual" é uma predição baseada no modelo para o cenário hipotético $\texttt{econ}=0$ nos respondentes com formação econômica; não expressa um experimento real, mas uma decomposição associativa sob as hipóteses do modelo (permutabilidade condicional não testável).

\begin{figure}[htbp]
\centering
\caption{Médias de concordância: "O déficit federal é grande demais".}
\includegraphics[width=0.80\textwidth]{Textuais/analise_v2/outputs/figures/counterfactual/déficit_federal_grande.png}
\notafigv{As três barras exibem: (1) média observada do grupo de controle (\texttt{econ}=0); (2) predição \textit{in-sample} do modelo para os economistas com \texttt{econ} zerado (decomposição condicional, não contrafactual causal — fatores não observados correlacionados com a escolha de fazer pós-graduação em Economia podem explicar parte ou toda a diferença visualizada); (3) média observada dos economistas (\texttt{econ}=1).}
\label{fig:deficit_v2}
\end{figure}

\begin{figure}[htbp]
\centering
\caption{Médias de concordância: "A produtividade está aumentando devagar demais".}
\includegraphics[width=0.80\textwidth]{Textuais/analise_v2/outputs/figures/counterfactual/produtividade_aumentando_devagar.png}
\notafigv{Neste item, o sinal se inverte: o grupo de economistas apresenta maior concordância com a preocupação de produtividade. A predição in-sample (barra central) quantifica a decomposição condicional associada ao preditor \texttt{econ}, condicionando nas covariáveis observadas; não implica que a formação seja a causa da diferença.}
\label{fig:produtividade_v2}
\end{figure}

\begin{figure}[htbp]
\centering
\caption{Médias de concordância: "Produtos importados são benéficos para a economia".}
\includegraphics[width=0.80\textwidth]{Textuais/analise_v2/outputs/figures/counterfactual/produtos_importados_benéficos.png}
\notafigv{O grupo de economistas exibe maior endosso à abertura comercial, consistente com o consenso documentado por \citeonline{Fuller2014}. A predição in-sample (barra central) é a decomposição condicional associada ao preditor \texttt{econ} — não um efeito causal da formação; auto-seleção e variáveis não observadas contribuem para a diferença entre os grupos.}
\label{fig:importados_v2}
\end{figure}

\paragraph{Déficit federal} (Figura~\ref{fig:deficit_v2}): o grupo de controle apresenta média de concordância $1{,}48$ com "o déficit é grande demais"; o grupo de economistas, $1{,}02$. A predição in-sample situa-se em $1{,}48$ — sugerindo, dentro das hipóteses do modelo, que a diferença está condicionalmente associada à formação econômica. A direção da diferença é consistente com o que \citeonline{Fuller2014} documentam empiricamente como posição predominante entre economistas: maior ceticismo com relação a alarmes fiscais de curto prazo, atenção à trajetória intertemporal da dívida e distinção entre déficit nominal e estrutural. Bastiat \cite{bastiat2010oqueseve} fornece um quadro intelectual para essa assimetria (o que se vê vs. o que não se vê), mas não deve ser interpretado como gabarito normativo: trata-se de referência literária que articula uma distinção analítica, não de âncora para definir qual resposta é "correta".

\paragraph{Produtividade} (Figura~\ref{fig:produtividade_v2}): o sinal se inverte. O grupo de controle apresenta média de concordância mais baixa com a preocupação de produtividade ($\bar{Y} = 1{,}12$); o grupo de economistas, mais alta ($\bar{Y} = 1{,}72$). A predição contrafactual ($1{,}12$) indica que essa diferença está associada à formação econômica, não a outras covariáveis. Isso é consistente com a hipótese de que o treinamento em econometria e macroeconomia torna a fricção de produtividade mais saliente — fricção que, sem instrumentais de mensuração da produtividade total dos fatores, o cidadão racional não tem incentivos para identificar \cite{downs1957economic}.

\paragraph{Produtos importados} (Figura~\ref{fig:importados_v2}): o grupo de economistas exibe média de endosso à abertura comercial mais alta ($\bar{Y} = 1{,}63$ vs $1{,}26$ do grupo de controle). A predição contrafactual ($1{,}26$) indica que a diferença está associada à formação econômica. Esse padrão é consistente com o consenso documentado por \citeonline{Fuller2014}: respondentes com formação econômica avançada tendem a posições mais favoráveis ao livre-comércio, alinhadas com a teoria das vantagens comparativas. Um padrão que Manzoni captou na ficção: diante de um preço que sobe ou de um mercado que parece hostil, o senso comum tende a preferir um culpado visível — o comerciante ganancioso, o estrangeiro oportunista — a um mecanismo impessoal de oferta e demanda.

\section{Efeitos Marginais Médios do Espectro Político}
\label{sec:ame-espectro}

O espectro político é o preditor dominante em termos de frequência (34/53 DVs BH-significativas) e de magnitude dos coeficientes. Para quantificar sua escala em termos interpretáveis de probabilidade, calculam-se AMEs via primeira diferença finita discreta em dois passos representativos da escala: Centro$\to$Direita ($s=0\to1$) e Esquerda$\to$Centro ($s=-1\to0$).

% AME do espectro político nos 10 itens de maior magnitude
% Gerado automaticamente por ame_espectro.py
\begin{table}[htbp]
\centering
\SingleSpacing
\caption{Efeitos Marginais Médios (AME) do Espectro Político nos 10 Itens de Maior Magnitude.}
\label{tab:ame_espectro_top10}
\footnotesize
\begin{tabular}{p{5.2cm}rrrrrr}
\toprule
 & \multicolumn{3}{c}{\textbf{Centro $\to$ Direita}} & \multicolumn{3}{c}{\textbf{Esq. $\to$ Centro}} \\
\cmidrule(lr){2-4}\cmidrule(lr){5-7}
\textbf{Item (DV)} & \textbf{$\Delta\bar{Y}$} & \textbf{AME\textsubscript{P=2}} & \textbf{$\hat{\beta}$} & \textbf{$\Delta\bar{Y}$} & \textbf{AME\textsubscript{P=2}} & \textbf{SE} \\
\midrule
Privatização de estatais é benéfica & 0,683 & 0,443 & 2,091*** & 0,620 & 0,158 & 0,271 \\
Governo atual sabe administrar a economia & -0,353 & -0,064 & -1,999*** & -0,675 & -0,265 & 0,293 \\
Governo regulamenta negócios demais & 0,529 & 0,377 & 1,842*** & 0,535 & 0,176 & 0,250 \\
Mulheres na força de trabalho causa desemprego & -0,352 & -0,302 & -1,427*** & -0,260 & -0,247 & 0,246 \\
Lucros empresariais à custa dos trabalhadores & -0,511 & -0,218 & -1,264*** & -0,490 & -0,295 & 0,198 \\
Brasil deveria adotar livre-comércio amplo & 0,288 & 0,254 & 1,145*** & 0,323 & 0,233 & 0,202 \\
Governo deve intervir nos preços & -0,328 & -0,093 & -1,079*** & -0,400 & -0,179 & 0,188 \\
Indústria nacional deve ser protegida & -0,380 & -0,199 & -1,048*** & -0,334 & -0,236 & 0,185 \\
Executivos ganham mais do que merecem & -0,343 & -0,108 & -1,075*** & -0,404 & -0,193 & 0,187 \\
O déficit federal é grande demais & 0,284 & 0,235 & 1,092*** & 0,349 & 0,234 & 0,198 \\
\bottomrule
\multicolumn{7}{c}{\footnotesize Fonte: elaboração do autor.} \\
\multicolumn{7}{p{14cm}}{\footnotesize \textit{Nota:} AME via primeira diferença finita discreta: $\Delta P(Y=j \mid \texttt{espectro}=s+1) - \Delta P(Y=j \mid \texttt{espectro}=s)$, médio sobre todos os $i$. Centro$\to$Direita: $s=0\to1$; Esq.$\to$Centro: $s=-1\to0$. $\Delta\bar{Y}$: variação na resposta esperada (escala 0--2). AME\textsubscript{P=2}: variação em P(concordo totalmente). $\hat{\beta}$: coeficiente logit ordenado. Estrelas: $p$-valor assintótico (*$p<0{,}05$; **$p<0{,}01$; ***$p<0{,}001$).} \\
\end{tabular}
\end{table}

A Tabela~\ref{tab:ame_espectro_top10} apresenta os resultados para os 10 itens com maior $|\hat{\beta}_{\text{esp}}|$. As magnitudes são expressivas: um passo na escala ideológica em direção à direita está associado a um aumento médio de $+0{,}68$ pontos (escala 0--2) na concordância com privatização de estatais e de $+0{,}53$ pontos na avaliação de regulação excessiva. Em sentido contrário, o mesmo passo está associado a $-0{,}51$ pontos na percepção de lucros excessivos e a $-0{,}35$ pontos na avaliação do governo atual. Essas diferenças são substancialmente maiores do que os efeitos condicionais da formação econômica nos quatro itens-chave (intervalo $|\Delta\bar{Y}| \in [0{,}14;\;0{,}60]$), confirmando quantitativamente que o espectro político é o preditor de maior magnitude nesta amostra.

Nota metodológica: os AMEs do espectro são baseados em modelos com $n\approx145$ observações válidas por item e não incluem intervalos de confiança bootstrap (o custo computacional com $B=999$ seria de $\approx 9.000$ estimações adicionais); os p-valores assintóticos reportados na tabela são suficientes para itens com $\hat{\beta}_{\text{esp}}$ da magnitude observada ($|\hat{\beta}|>1{,}0$, SE~$<0{,}3$).

\section{Diagnósticos e Robustez}
\label{sec:diagnosticos-resultados}

\subsection{Pressupostos dos modelos}

O teste de \citeonline{Brant1990} confirmou a validade do pressuposto de proporcionalidade das chances em todas as 48 DVs estimadas por PO: nenhuma falha ($p_{\min} > 0{,}32$). O VIF máximo de $2{,}1$ confirma ausência de multicolinearidade problemática. Resultados completos em \autoref{apendice:tab_brant_v2}.

\subsection{Intervalos bootstrap e testes de permutação}

Para os quatro itens-chave, os intervalos BC$_0$ (bootstrap, $B=999$) são detalhados em \autoref{apendice:tab_bootstrap_v2}. Os ICs para \textit{déficit federal} e \textit{produtividade} excluem zero, reforçando os efeitos nominais mesmo sob quantificação de incerteza robusta. Os testes de permutação ($P = 5.000$) confirmam: p-valores abaixo de $0{,}05$ para esses dois itens (\autoref{apendice:tab_permutation_v2}).

\subsection{McFadden $R^2$ e ajuste geral}

A mediana de McFadden $R^2 = 0{,}099$ indica ajuste modesto mas adequado para variáveis de opinião política — em linha com a literatura de modelos de crenças políticas \cite{mccullaghNelder1989}. Distribuição completa em \autoref{apendice:tab_gof_v2}.

\section{Síntese Frente às Hipóteses H1--H7}
\label{sec:sintese-hipoteses}

\textbf{H1 — Vieses populares sistemáticos:} \textit{Corroborada}. Padrões compatíveis com vieses antimercado, antiestrangeiro, antitrabalho e pessimista são observados em múltiplos itens e blocos, persistentes sob controle de covariáveis.

\textbf{H2 — Conhecimento reduz vieses:} \textit{Parcialmente corroborada, sujeita a limitação de poder}. Cinco efeitos nominais com direção correta; nenhum sobrevive a BH. A análise de poder ($28{,}2\%$) indica que o resultado nulo é previsível mesmo na presença de associações reais moderadas. Os itens de livre-comércio e produtividade oferecem a diferença condicional mais clara, consistente com a hipótese de que formação econômica está associada a respostas distintas em itens com componente factual dominante.

\textbf{H3 — Viés antimercado:} \textit{Corroborada}. Taxa BH-espectro de 9/9 (100\%), confirmando clivagem sistemática entre esquerda e direita nas percepções de mercado. A formação econômica não tem efeito BH-significativo neste bloco, sugerindo que as clivagens são predominantemente valorativas e identitárias.

\textbf{H4 — Viés antiestrangeiro:} \textit{Parcialmente corroborada}. Bloco H4 apresenta 5/9 DVs BH-espectro-significativas e \textit{produtos importados benéficos} emerge como o item com maior diferença condicional associada à formação ($\Delta\bar{Y} = +0{,}37$, IC-bootstrap excluindo zero). A diferença associada à formação aparece em itens de comércio; não em imigração.

\textbf{H5 — Viés antitrabalho:} \textit{Parcialmente corroborada}. Bloco H5 com 6/10 DVs BH-espectro-significativas. A maior preocupação de economistas com produtividade sugere que, em itens de fricção estrutural, a formação funciona como amplificador de atenção legítima, não como mitigador de alarme infundado.

\textbf{H6 — Viés pessimista:} \textit{Corroborada}. Bloco H6 com 4/10 DVs BH-espectro-significativas. Economistas são mais preocupados com produtividade estrutural do que o público leigo — inversão de sinal que revela a assimetria informacional.

\textbf{H7 — Ideologia como filtro:} \textit{Fortemente corroborada}. Bloco H7 apresenta 11/12 DVs BH-espectro-significativas — a maior taxa do estudo. Reformas, avaliação do governo, taxa Selic e corrupção são percebidas através da lente ideológica, com espectro explicando a maior parcela da variância em cada item. Este é o resultado mais robusto do estudo.

\begin{figure}[htbp]
\centering
\caption{Volcano plot: coeficiente de formação econômica vs.\ $-\log_{10}(p)$ para as 53 DVs.}
\includegraphics[width=0.90\textwidth]{Textuais/analise_v2/outputs/figures/forestplots/volcano_econ.png}
\notafigv{Eixo horizontal: $\hat{\beta}_{\texttt{econ}}$ (positivo = economistas concordam mais). Eixo vertical: $-\log_{10}(p)$; linha tracejada = $p = 0{,}05$. A dispersão horizontal revela a diversidade de direções; a ausência de pontos acima do limiar BH evidencia o problema de poder, não a ausência de padrão.}
\label{fig:volcano_econ}
\end{figure}

O volcano plot (Figura~\ref{fig:volcano_econ}) encapsula o diagnóstico: efeitos nominais dispersos em ambas as direções, concentrados em magnitude pequena a moderada, nenhum sobrevivendo ao limiar BH. A distribuição sugere efeitos reais de magnitude moderada ($|\beta| \approx 0{,}5$--$0{,}9$, equivalente a OR $\approx 1{,}6$--$2{,}5$), detectáveis com $n_{\texttt{econ}} \geq 60$--$120$.

A tabela síntese completa dos 53 DVs consta de \autoref{apendice:tab_sintese_v2}.
